% GENERATED FILE. Edit canonical lessons, then run npm run generate:books.
% canonical-source-hash: fnv1a64:ee9fd8c39ce162b8
% canonical-lessons: ES-C701-caluroso, ES-C701-templado, ES-C701-bochorno, ES-C701-sequia, ES-C701-sudar

\chapter{Calor y sequía}
\label{ch:es-calor-y-sequia}
\hlchaptermodality{701}

\begin{chapteropening}
\textbf{By the end of this chapter:} \emph{I can talk about warm weather: mild, hot and muggy days, a drought, and sweating.}

Complete the last lesson of chapter 701: I can talk about warm weather: mild, hot and muggy days, a drought, and sweating.
\end{chapteropening}

\section[caluroso]{caluroso — hot (weather); warm (welcome)}

\label{lesson:ES-C701-caluroso}



\begin{quote}
Before the new one: say the Spanish for the snowfall, then the Spanish for the hoarfrost.
\end{quote}

\subsection*{You'll want to know: caluroso}
\textbf{caluroso} — \textquotedblleft{}hot (weather); warm (welcome)\textquotedblright{}.

\begin{quote}
Fue un verano muy caluroso.
\end{quote}

\begin{quote}
\emph{It was a very hot summer.}
\end{quote}

\emph{Un caluroso saludo} — a warm greeting.

\begin{cousinweb}
\emph{caluroso} is built on \emph{calor} \textquotedblleft{}heat\textquotedblright{}, which is inherited from Latin \emph{calor} (EtymDB).
\end{cousinweb}

\begin{grammarlens}[title={What you've built}]
One more word for telling, explaining and describing at B1.
\end{grammarlens}

\subsection*{Your turn}
\begin{itemize}
\raggedright
  \item \emph{Say it:} \emph{caluroso}
  \item \emph{Say it:} \emph{caluroso}, once more
  \item \emph{Say it:} a sentence of your own with \emph{caluroso}
  \item \emph{Recall:} say the Spanish for the snowfall, then the Spanish for the hoarfrost, then say \emph{caluroso} again
\end{itemize}

\begin{tcolorbox}[breakable,colback=teal!4,colframe=teal!35!black,title={Before you move on}]
What does caluroso mean? (\textquotedblleft{}Hot (weather); warm (welcome)\textquotedblright{}.) Say it once more.
\end{tcolorbox}
\section[templado]{templado — mild, temperate; lukewarm}

\label{lesson:ES-C701-templado}



\begin{quote}
Before the new one: say the Spanish for the hoarfrost, then the Spanish for hot.
\end{quote}

\subsection*{You'll want to know: templado}
\textbf{templado} — \textquotedblleft{}mild, temperate; lukewarm\textquotedblright{}.

\begin{quote}
En primavera el tiempo es templado.
\end{quote}

\begin{quote}
\emph{In spring the weather is mild.}
\end{quote}

\emph{Agua templada} — lukewarm water.

\begin{cousinweb}
\emph{templado} is the past form of \emph{templar} \textquotedblleft{}to temper\textquotedblright{} (EtymDB), Old Spanish \emph{temprar}, from Latin \emph{temperāre} \textquotedblleft{}to mix in due measure\textquotedblright{}.
\end{cousinweb}

\begin{grammarlens}[title={What you've built}]
One more word for telling, explaining and describing at B1.
\end{grammarlens}

\subsection*{Your turn}
\begin{itemize}
\raggedright
  \item \emph{Say it:} \emph{templado}
  \item \emph{Say it:} \emph{templado}, once more
  \item \emph{Say it:} a sentence of your own with \emph{templado}
  \item \emph{Recall:} say the Spanish for the hoarfrost, then the Spanish for hot, then say \emph{templado} again
\end{itemize}

\begin{tcolorbox}[breakable,colback=teal!4,colframe=teal!35!black,title={Before you move on}]
What does templado mean? (\textquotedblleft{}Mild, temperate; lukewarm\textquotedblright{}.) Say it once more.
\end{tcolorbox}
\section[el bochorno]{el bochorno — the muggy heat, close weather; embarrassment}

\label{lesson:ES-C701-bochorno}



\begin{quote}
Before the new one: say the Spanish for hot, then the Spanish for mild.
\end{quote}

\subsection*{You'll want to know: el bochorno}
\textbf{el bochorno} — \textquotedblleft{}the muggy heat, close weather; embarrassment\textquotedblright{}.

\begin{quote}
¡Qué bochorno hace hoy!
\end{quote}

\begin{quote}
\emph{It is so muggy today!}
\end{quote}

It also means embarrassment: \emph{¡qué bochorno!}

\begin{cousinweb}
\emph{bochorno} is inherited from Latin \emph{vulturnus} \textquotedblleft{}east wind\textquotedblright{} (EtymDB); the name of a hot wind became the name for hot, heavy weather.
\end{cousinweb}

\begin{grammarlens}[title={What you've built}]
One more word for telling, explaining and describing at B1.
\end{grammarlens}

\subsection*{Your turn}
\begin{itemize}
\raggedright
  \item \emph{Say it:} \emph{el bochorno}
  \item \emph{Say it:} \emph{el bochorno}, once more
  \item \emph{Say it:} a sentence of your own with \emph{el bochorno}
  \item \emph{Recall:} say the Spanish for hot, then the Spanish for mild, then say \emph{el bochorno} again
\end{itemize}

\begin{tcolorbox}[breakable,colback=teal!4,colframe=teal!35!black,title={Before you move on}]
What does el bochorno mean? (\textquotedblleft{}The muggy heat, close weather; embarrassment\textquotedblright{}.) Say it once more.
\end{tcolorbox}
\section[la sequía]{la sequía — the drought}

\label{lesson:ES-C701-sequia}



\begin{quote}
Before the new one: say the Spanish for mild, then the Spanish for the muggy heat.
\end{quote}

\subsection*{You'll want to know: la sequía}
\textbf{la sequía} — \textquotedblleft{}the drought\textquotedblright{}.

\begin{quote}
Este año hay sequía y no llueve.
\end{quote}

\begin{quote}
\emph{There is a drought this year and it does not rain.}
\end{quote}

Built on \emph{seco}, dry.

\begin{cousinweb}
\emph{sequía} is built on \emph{seco} \textquotedblleft{}dry\textquotedblright{} (EtymDB), which is inherited from Latin \emph{siccus} \textquotedblleft{}dry\textquotedblright{}.
\end{cousinweb}

\begin{grammarlens}[title={What you've built}]
One more word for telling, explaining and describing at B1.
\end{grammarlens}

\subsection*{Your turn}
\begin{itemize}
\raggedright
  \item \emph{Say it:} \emph{la sequía}
  \item \emph{Say it:} \emph{la sequía}, once more
  \item \emph{Say it:} a sentence of your own with \emph{la sequía}
  \item \emph{Recall:} say the Spanish for mild, then the Spanish for the muggy heat, then say \emph{la sequía} again
\end{itemize}

\begin{tcolorbox}[breakable,colback=teal!4,colframe=teal!35!black,title={Before you move on}]
What does la sequía mean? (\textquotedblleft{}The drought\textquotedblright{}.) Say it once more.
\end{tcolorbox}
\section[sudar]{sudar — to sweat}

\label{lesson:ES-C701-sudar}



\begin{quote}
Before the new one: say the Spanish for the muggy heat, then the Spanish for the drought.
\end{quote}

\subsection*{You'll want to know: sudar}
\textbf{sudar} — \textquotedblleft{}to sweat\textquotedblright{}.

\begin{quote}
Con este calor sudo mucho.
\end{quote}

\begin{quote}
\emph{I sweat a lot in this heat.}
\end{quote}

\emph{El sudor} is sweat.

\begin{cousinweb}
\emph{sudar} is inherited from Latin \emph{sūdāre} \textquotedblleft{}to sweat\textquotedblright{} (EtymDB).
\end{cousinweb}

\begin{grammarlens}[title={What you've built}]
One more word for telling, explaining and describing at B1.
\end{grammarlens}

\subsection*{Your turn}
\begin{itemize}
\raggedright
  \item \emph{Say it:} \emph{sudar}
  \item \emph{Say it:} \emph{sudar}, once more
  \item \emph{Say it:} a sentence of your own with \emph{sudar}
  \item \emph{Recall:} say the Spanish for the muggy heat, then the Spanish for the drought, then say \emph{sudar} again
\end{itemize}

\begin{tcolorbox}[breakable,colback=teal!4,colframe=teal!35!black,title={Before you move on}]
What does sudar mean? (\textquotedblleft{}To sweat\textquotedblright{}.) Say it once more.
\end{tcolorbox}
